% !TeX root = ../../../navier-stokes-manuscript.tex
% ============================================================
% 2.4 The Spatial Column Inside Critical Height
% ============================================================

\subsection{The Spatial Column Inside Critical Height}\label{sub:TheSpatialColumnInsideCriticalHeight}

For every real \(s\geq0\), define the spatial Sobolev energy
\[
  E_s(t)
  :=
  \frac12\norm{\Lambda^su(t)}_2^2.
\]
The critical height from the scaling calculation is the particular row \(H=E_{1/2}\).
The value of introducing all \(E_s\) at once is that viscosity moves them by one complete Sobolev step.

Pair the momentum equation with \(\Lambda^{2s}u\).
Equivalently, apply \(\Lambda^s\) and take the \(L^2\) inner product with \(\Lambda^su\).
The time derivative gives \(E_s'\), and integration by parts gives
\[
  \nu\pair{\Lambda^su}{\Lambda^s\Delta u}
  =
  -\nu\norm{\Lambda^{s+1}u}_2^2
  =
  -2\nu E_{s+1}.
\]
Keep the transport and pressure response together in the signed production term
\[
  \mathcal P_s(t)
  :=
  -\pair{\Lambda^su}
  {\Lambda^s\bigl((u\cdot\nabla)u+\nabla p\bigr)}.
\]
For a smooth decaying divergence-free field, the global pressure pairing vanishes because
\[
  \pair{\Lambda^su}{\Lambda^s\nabla p}
  =
  -\pair{\nabla\cdot\Lambda^su}{\Lambda^sp}
  =0.
\]
We retain pressure in the definition because the balance was derived from the pressure-complete equation and every local differentiated row still contains its pressure coordinate.
The exact Sobolev balance is
\[
  E_s'
  =
  \mathcal P_s-2\nu E_{s+1}.
\]
The production term records the signed nonlinear transfer into height \(s\), while viscosity exposes the next spatial height \(s+1\).

At the critical row, the first time derivative is
\[
  H'
  =
  \mathcal P_{1/2}-2\nu E_{3/2}.
\]
One temporal change in critical height has already opened a spatial energy one full Sobolev step higher.
The row also retains the signed production needed to distinguish transfer into that height from viscous removal at the next height.

Differentiate once more.
Before compressing the result, use the same energy balance at \(s=3/2\):
\[
  E_{3/2}'
  =
  \mathcal P_{3/2}-2\nu E_{5/2}.
\]
Substitution gives
\begin{align*}
  H''
  &={}
  \mathcal P_{1/2}'-2\nu E_{3/2}'
  \\
  &={}
  \mathcal P_{1/2}'
  -2\nu\mathcal P_{3/2}
  +(2\nu)^2E_{5/2}.
\end{align*}
The second temporal derivative exposes \(E_{5/2}\), while the lower spatial heights survive through the differentiated production terms.

A third differentiation uses
\[
  E_{5/2}'
  =
  \mathcal P_{5/2}-2\nu E_{7/2}.
\]
The first three critical-height rows are
\begin{align*}
  H'
  &={}
  \mathcal P_{1/2}-2\nu E_{3/2},
  \\
  H''
  &={}
  \mathcal P_{1/2}'
  -2\nu\mathcal P_{3/2}
  +(2\nu)^2E_{5/2},
  \\
  H'''
  &={}
  \mathcal P_{1/2}''
  -2\nu\mathcal P_{3/2}'
  +(2\nu)^2\mathcal P_{5/2}
  -(2\nu)^3E_{7/2}.
\end{align*}
The alternating powers of \(2\nu\) come from repeatedly substituting the viscous part of the next spatial energy balance.
Each substitution also leaves behind the production term at the height it passes.

The pattern can now be stated without hiding any lower row.
For every integer \(n\geq1\), repeated differentiation and substitution give
\[
  H^{(n)}
  =
  (-2\nu)^nE_{n+1/2}
  +
  \sum_{j=0}^{n-1}
  (-2\nu)^j
  \partial_t^{\,n-1-j}\mathcal P_{j+1/2}.
\]
The formula follows by induction: differentiating the \(n\)-th row differentiates every existing production term, and replacing \(E_{n+1/2}'\) with \(\mathcal P_{n+1/2}-2\nu E_{n+3/2}\) creates exactly the new final summand and the next highest energy.
No production derivative is discarded in the passage to higher order.

This triangular identity is the scalar shadow of the pressure-complete mixed jet.
It shows that temporal response at the critical height is nested with the spatial Sobolev column
\[
  E_{3/2},
  \quad
  E_{5/2},
  \quad
  E_{7/2},
  \quad\ldots
\]
through all of the signed lower production responses that the equation generates.
Temporal order and spatial height remain different coordinates.
The identity relates them; it does not collapse them into a single derivative count.

The identity also carries no boundedness conclusion by itself.
Large production and large viscous height can cancel in one row, and their derivatives can interact differently in the next.
Its payment is structural: a critical-height failure cannot be analysed honestly while ignoring the spatial energies and lower production derivatives nested inside it.
Once every finite row is retained, Sobolev embedding tells us exactly what their simultaneous all-orders control would mean.
